% Drawing 138/06-42 -- folder 138, batch 06, archive page 117 (PDF page 118).
% Transcription: un segment suivi d'une fourche, marqué $\times$ ; un segment suivi d'une boucle, marqué $\times$ ; un cercle portant deux sommets avec une arête vers l'extérieur.
% Kind: small maps (degree 5). Method 2 (schematic redraw from the facsimile). Status: ok.
% Notes: c): a segment followed by a fork, marked ×; a segment followed by a loop, marked ×; a circle carrying two vertices with an edge going outwards. His small crosses beside some maps are kept as × marks. Small planar maps listed by degree (number of edges) in the text; redrawn, not encoded (each would be a map/1 with at most a few edges; the text gives the classification).
% Coordinates: 1 unit = 1 mm = 1/100 of the crop width (bbox in page fractions: [0.632, 0.447, 0.847, 0.631]).
% Redrawn by Claude (claude-opus-5-5), issue #73, 2026-11-07; not checked by a human; not compiled (no TeX in the session).
% Rights: reproduces part of the work -- behind RENDER_TRANSCRIPTS, never published by itself.
\documentclass[tikz,border=3pt]{standalone}
\usepackage{amsmath,amssymb}
\usetikzlibrary{patterns,arrows.meta}
\tikzset{pencil/.style={black!45}, struck/.style={black!35}, mark/.style={circle,draw,inner sep=1.3pt}, pt/.style={circle,fill,inner sep=0.9pt}, opt/.style={circle,draw,fill=white,inner sep=1.1pt}}
\begin{document}
\begin{tikzpicture}[x=1mm,y=1mm,line cap=round,line join=round,font=\scriptsize]
\draw (4.6,-15.4) -- (46.2,-15.4);
\draw (46.2,-15.4) -- (60,-6.2);
\draw (46.2,-15.4) -- (60,-26.2);
\node[pt] at (24.6,-16.9) {};
\node[pt] at (46.2,-15.4) {};
\node[inner sep=1pt] at (69.2,-15.4) {\tiny $\times$};
\draw (6.2,-43.1) -- (46.2,-41.5);
\draw (60.8,-43.8) ellipse[x radius=14.6, y radius=10];
\node[pt] at (24.6,-43.1) {};
\node[pt] at (46.2,-41.5) {};
\node[inner sep=1pt] at (86.2,-49.2) {\tiny $\times$};
\draw (34.6,-93.8) ellipse[x radius=22.3, y radius=21.5];
\node[pt] at (12.3,-92.3) {};
\node[pt] at (56.9,-92.3) {};
\draw (56.9,-92.3) -- (87.7,-92.3);
\end{tikzpicture}
\end{document}
